\documentclass[11pt]{article}
\usepackage{amsmath,hyperref}
\title{Can a fleet commit to future entry without granting current entry?}
\date{25 November 2026}
\begin{document}
\maketitle
\paragraph{Live duty and authority.}
A support crew can spend \(c\) nonrefundable work credits before a rescue route opens. The old users occupy gates owned by separate voluntary operators. The crew may enter only after each gate's actual old duty settles and a local fence makes reentry by the old job impossible. A receipt must reach the part of the system that can make the entry effective, not merely be signed or displayed. A refused commercial or operational consent cannot relax the physical constraint. Contextual timing follows the existing \href{https://github.com/eternal-november/eos-research-programmes/blob/5003ee534aadd3614868022c93961029d19333e3/projects/byzantine-control/portability/HANDOFF-RECEIPT-REUSE.md}{receipt-reuse study}: retain local readiness, actual delivery and competing use on the same history. This fixture assumes authenticated correct physical release and otherwise refuses. It does not produce such evidence from inconsistent reports.
\paragraph{Timing consequence.}
For each gate, \(t_i\) is the time qualified release evidence would support effective entry, including its propagation; absent evidence means \(t_i=\infty\). A gate with no existing commitment imposes an additional approval delay \(\ell_i\), taken here as a separately available effective delay, even granting simultaneous requests as the optimistic comparison. A binding instruction signed earlier may vest the entry role only at \(t_i\). For travel duration \(d\) and deadline \(D\), the exact condition is
\[
 \max_i t_i + d\le D
\]
for prior conditional consent, and \(\max_i(t_i+\ell_i)+d\le D\) when every consent is only sought after release. Later release still prevents service in both. Arrival of the old receipt alone can never replace consent. This construction uses local gate enforcement and a crew that can gather all required readiness; it does not suppose that all gate owners confer at release time. Restarts, expired grants, conflicting gate schedules, and shared travel alternatives need their own source law.
\paragraph{Voluntary participation.}
In a separate price subgame, let \(b\) be the gross value of timely completed rescue to the crew. Gate owner \(i\) loses a known net outside value \(r_i\ge0\) if the corresponding role is used. An authorized toll \(p_i\) is credited only on actual, timely completed service, is never debited on an unsuccessful attempt, and cannot exceed the gross budget in aggregate. No rescuer or owner can buy a waiver of physical release. Work and tolls are illustrative integers, not a payment implementation. Physical releases and deadline feasibility are fixed in this subgame; failure yields owners their outside alternatives and gives the crew no service return.
\par
When the right to name a toll is retained until after preparation, owners announce tolls in sequence and any may refuse. Earlier announcements cannot be repriced within that subgame; the last gate can unilaterally set its toll. The crew sees all final tolls and chooses whether to use the authorized route, with cost \(c\) already sunk. At any accepted pure subgame-perfect price outcome, at most one credit remains to the crew after tolls. If two or more remained, the last owner could raise its toll by one and the crew would still strictly prefer entry to refusal, so such an outcome cannot be optimal for the last owner. This argument does not assume a favorable tie resolution. If \(c\ge2\), the crew does not prepare, even in a case where a coordinator could choose a mutually beneficial full allocation. The obstruction presumes there is no earlier enforceable fee agreement or effective toll cap. It is a version of the familiar hold-up problem, not a novel equilibrium concept.
\par
For this fixed-success branch, with a binding pre-release option, each owner can agree to its own future toll and vest the crew's local right on qualified physical release. The owner can decline before agreeing, and one that did agree cannot revoke after the preparation cost becomes sunk under the stipulated interface. Strict participation for all owners and the crew is feasible exactly when
\[
 b\ge c+\sum_i r_i + m+1 ,
\]
where \(m\) is the number of owners. Each integer toll must exceed \(r_i\) by at least one and one credit must remain to the crew after preparation; choosing \(p_i=r_i+1\) attains the bound. With weak participation remove \(m+1\), but then a preferred tie resolution is not guaranteed. No option under this model grants present entry or makes its other branch safe. Suppose \(r=(1,2), b=8,c=2\), and the gate evidence becomes usable at ticks 2 and 3, each post-release approval would add two ticks, and travel takes two ticks by deadline 5. Prices 2 and 3 give both owners and the crew a net gain of one, with timely service under the option. Without enforceable early consent costly preparation fails even ignoring approval delays; waiting adds a second independent obstruction.

\paragraph{Uncertain release is one joint event.}
Suppose now that before preparing, all parties face a fixed declared probability \(q\) that every qualifying receipt arrives on time and the authorized rescue can complete. Call that the success event; all tolls and the \(r_i\) losses occur only on it. Let gate owner \(i\) additionally incur an upfront cost \(k_i\geq0\) on committing. With risk-neutral values on this declared law, strict participation is exactly
\[
 q>0,\qquad b\geq \sum_i \bigl(r_i+\lfloor k_i/q\rfloor+1\bigr)
          +\lfloor c/q\rfloor+1 .
\]
Indeed each owner needs \(q(p_i-r_i)>k_i\), so use the smallest permitted positive margin, and the crew needs \(q(b-\sum_i p_i)>c\). Capacity and authority still have to permit actual service on every counted success history. This notional welfare calculation never licenses relaxation of a safety duty. With two gates and each timely with probability \(7/8\), their joint probability can be \(7/8\) when losses coincide or \(3/4\) when they are disjoint. With \(r=(1,2), b=9,c=3,k=(0,0)\), the same minimum tolls give strict preparation in the first law and only a tie in the second. No independent-resampling assumption or field probability is being made. The rational replay checks the strict fee bound directly on several such joint success rates.
\paragraph{Baselines and checks.}
The administrative baseline is \href{https://doi.org/10.1145/300830.300839}{Sandhu, Bhamidipati and Munawer's ARBAC97} user-role assignment with prerequisites. If the actual clearance is represented by a reliable local membership update, a correctly specified direct grant guard can implement the same conditional vesting. ARBAC assignment by itself neither supplies the evidence nor makes an economic commitment binding. \href{https://pages.nyu.edu/debraj/Papers/KonishiRay.pdf}{Konishi and Ray} treat coalition change as an ongoing process with endogenous continuation values. Here only an ordinary finite, fixed-trajectory subgame is solved; the model is not their general coalition equilibrium. An equal-interface extensive-form comparison agrees, so there is no advantage over a direct scheme with commitment, right binding and physical witnesses. Pure-strategy equilibrium outcomes for one to three owners were exhaustively enumerated with arbitrary tie choices. Separately the timing replay searches all permitted start ticks, including absent evidence. Lean checks the integer-allocation and last-price algebra; it does not formalize the whole game. This is authored partial R5, not evidence for Byzantine reports or an independent transfer.
\end{document}
